\subsection{Countable simultaneous systems that have no solution}

\begin{claim}{A.27}{Proved}{Explicit diagonal systems with identity exponent matrix}
Let $E$ be a nontrivial countable group satisfying at least one of
the following three properties:
\begin{enumerate}
\item every nonidentity word in $E*\langle t\rangle$ has a
nonidentity evaluation at some $t\in E$;
\item $E$ is torsion-free and has exactly two conjugacy classes;
\item there is an integer $N\geq1$ such that, for every
$a\ne1$ in $E$, every element of $E$ is a product of at most
$N$ conjugates of $a$ or $a^{-1}$.
\end{enumerate}
In case~3 assume also that $E$ is nonabelian. Then there is a
specified countable system of equations over $E$, with one
equation for each variable and exponent-sum matrix exactly the
identity, that has no solution in $E$.

More generally, let $H$ be torsion-free, let $K\leq H$ be
any subgroup, let $q:K\twoheadrightarrow E$ be an epimorphism,
and let $B$ be a specified CW classifying space for $H$.
There is a classifying space $Y$
containing $B$, obtained by countably many one- and two-cell
attachments, such that the actual inclusion gives
\[
 H\hookrightarrow P=\pi_1(Y),\qquad
 H_*(B;\mathbb Z)\xrightarrow{\cong}H_*(Y;\mathbb Z),
 \qquad \dim Y\leq\max(\dim B,2).
 \tag{A.27.1}
\]
The image of $H$ normally generates $P$, and no homomorphism
$P\to E$ restricts to $q$ on $K$. No extension of $q$ to $H$
is assumed. In particular a specified two-dimensional
model for $H$ gives an actual two-dimensional model for $P$.

Every nontrivial countable group existentially closed among all
groups has property~1 and consequently fails to solve all
countable systems with integral invertible exponent matrix.
The same conclusion holds for a nontrivial countable group
existentially closed among torsion-free groups. Moreover the
property of solving every such countable system passes to
epimorphic images. A group with that property therefore cannot
have a nontrivial countable quotient in any of the three
classes above.
\end{claim}

\begin{proof}
We use $[a,b]=aba^{-1}b^{-1}$ throughout. First suppose $E$
has property~1. This property implies that $E$ is infinite:
the word $t^m$, for a common multiple $m$ of all element orders,
would be a nonidentity mixed identity in any finite group.
Choose $h\ne1$, choose $b_0\in E$, and enumerate $E$ without
repetitions as $e_1,e_2,\ldots$. We shall recursively choose
$b_1,b_2,\ldots$ so that the equations
\[
 \begin{split}
 x_0&=[h,x_1]b_0,\\
 x_n&=[x_0e_n^{-1},x_{n+1}]b_n\qquad(n\geq1)
 \end{split}
 \tag{A.27.2}
\]
have no solution in $E$.

Suppose the coefficients with indices less than $k$ have been
chosen. Set $p=e_k$ and form words in $E*\langle t\rangle$
\[
 f_0(t)=[h,t]b_0,\qquad
 f_j^p(t)=[pe_j^{-1},t]b_j\quad(1\leq j<k),
 \qquad
 F_k^p=f_0\circ f_1^p\circ\cdots\circ f_{k-1}^p.
 \tag{A.27.3}
\]
Every coefficient $pe_j^{-1}$ appearing here is nonidentity.
Free-product reduced normal form makes the free factor $E$
malnormal in $E*\langle t\rangle$: if $w\notin E$, then
$E\cap wEw^{-1}=\{1\}$. Consequently, for $a\in E-\{1\}$
and $w\notin E$, the word $[a,w]$ does not belong to $E$.
Otherwise the equality
\[
 wa^{-1}w^{-1}=a^{-1}[a,w]
\]
would put a nonidentity element in that intersection.
Multiplication on the right by an element of $E$ keeps a word
outside $E$. Starting with $t$ and applying this observation
from the innermost substitution outwards gives $F_k^p(t)\notin E$.
Thus $F_k^p(t)p^{-1}$ is a nonidentity word. Property~1 supplies
$b_k\in E$ with
\[
 F_k^{e_k}(b_k)\ne e_k.
 \tag{A.27.4}
\]
Fixing an enumeration and always taking its first suitable
element specifies the recursion completely.

In a solution of (A.27.2), the value of $x_0$ would be some
$e_k$. Its $k$th equation would give $x_k=b_k$. Substitution
through the preceding equations would then yield
$x_0=F_k^{e_k}(b_k)\ne e_k$, a contradiction.

Property~2 gives the same choices by a separate direct argument.
For $a\ne1$ the map
\[
 f_{a,b}:E\longrightarrow E,\qquad y\longmapsto[a,y]b
 \tag{A.27.5}
\]
has image $E-\{ab\}$, because
$[a,y]b=a(ya^{-1}y^{-1})b$ and the middle factor runs through
all nonidentity elements. Each nonempty fiber is a right coset
of $C_E(a)$ and is infinite: this centralizer contains the
infinite cyclic subgroup generated by $a$. Restricting a map
with infinite nonempty fibers to a cofinite subset does not
change its image. Induction on the composition in (A.27.3)
therefore gives
\[
 F_k^{e_k}(E)=E-\{hb_0\}.
 \tag{A.27.6}
\]
This image is infinite, so (A.27.4) can again be enforced.
No assertion that every group in property~2 has property~1 is
needed for this argument.

We next treat property~3. It implies that every nontrivial
normal subgroup is all of $E$; thus $E$ is nonabelian simple
and has trivial center. Put $d=2N$, and define the word
\[
 W_a(u_1,v_1,\ldots,u_N,v_N)
     =\prod_{i=1}^{N}v_i[a,u_i]v_i^{-1}.
 \tag{A.27.7}
\]
It is identically $1$ when $a=1$ and has exponent sum zero
in every displayed variable and in $a$. For $a\ne1$, choose
$u$ with $c=[a,u]\ne1$. Property~3 expresses each element as
a product of at most $N$ conjugates of $c$ and $c^{-1}$.
The identity
\[
 c^{-1}=u[a,u^{-1}]u^{-1}
\]
puts either sign in the form of a conjugate of $[a,u_i]$.
Unused factors can be made $1$ by setting $u_i=1$.
Consequently the map $W_a:E^d\to E$ is surjective whenever
$a\ne1$.

For finite $E$ first enumerate its elements without repetition
as $e_1,\ldots,e_m$ and then extend this list periodically;
for infinite $E$ use the countable repetition-free enumeration
above. Use variables $x_s$ at all vertices of the infinite
rooted $d$-ary tree: $s$ ranges over the finite words in
$\{1,\ldots,d\}$. Write $o$ for the root. All variables at
depth $n\geq1$ have equations with the same coefficient
$x_oe_n^{-1}$. The equations are
\[
 \begin{aligned}
 x_o&=W_h(x_1,\ldots,x_d)b_o,\\
 x_s&=W_{x_oe_n^{-1}}(x_{s1},\ldots,x_{sd})b_s
       &&(|s|=n\geq1).
 \end{aligned}
 \tag{A.27.8}
\]
Only the first $m$ levels will enter the contradiction when
$E$ is finite. Every variable in either case has exactly one
equation, and the variable and equation sets are countable.

Here are the choices of the coefficients. Choose $b_o$ first.
After the choices at depths less than $n$, substitute
$x_o=e_n$ in the right sides of those earlier equations.
All coefficients $e_ne_j^{-1}$ with $j<n$ are nonidentity,
provided $n$ is in the first, repetition-free enumeration.
Each map at these levels is therefore surjective, including
the root map with coefficient $h$. Because each vertex has
its own disjoint list of children, successive substitution
defines a surjective map
\[
 T_n:E^{d^n}\twoheadrightarrow E
 \tag{A.27.9}
\]
from the values at depth $n$ to the computed value at the root.
Choose the finitely many coefficients $b_s$, $|s|=n$, so that
$T_n((b_s))\ne e_n$. Such a choice exists since $E\ne1$.
In the finite case all coefficients beyond depth $m$ may be
chosen arbitrarily. In any solution, $x_o=e_n$ for one of
the repetition-free indices. Then every equation at depth
$n$ has $W_1=1$, forcing $x_s=b_s$ there. The earlier
equations force $x_o=T_n((b_s))\ne e_n$, the required
contradiction.

In (A.27.2) the relator with index $n$ has exponent column
exactly the basis vector for $x_n$. In (A.27.8) the same is
true of the relator indexed by $s$: (A.27.7) has zero
exponents both in its children and in its coefficient
$x_oe_n^{-1}$. Thus both constructions have actual relative
boundary the identity on the indicated free abelian direct
sums, not merely a matrix that becomes invertible over
$\mathbb Q$.

We now construct the spaces in (A.27.1). Lift every chosen
coefficient through $q:K\twoheadrightarrow E$ and regard its
chosen lift as an element of the specified subgroup $K\leq H$,
using the same
lift whenever that coefficient occurs in the formulas.
Attach all the new variable loops to $B$. In the chain
construction, let $Y_N$ include the two-cells of indices
$0,\ldots,N$. Eliminate the defining generator--relator pairs
in the order $x_N,x_{N-1},\ldots,x_1$, substituting each
expression into the earlier relators. These are the usual
presentation-complex homotopy equivalences, relative to $B$
and to every loop not eliminated. The result is the wedge of
$B$ with the tail circles $x_{N+1},x_{N+2},\ldots$, with
one new loop $x_0$ and one two-cell. Its remaining relator
has total $x_0$-exponent $1$: all occurrences on the right
side lie inside the outer commutator with the lift of $h$.

For the tree construction, let $Y_N$ include the two-cells
at depths at most $N$. Eliminate all nonroot variables at
these depths, in descending order of depth. Each defining
right side uses only the root and children at the next depth,
so it does not contain the variable being eliminated.
Again the result is $B$ wedged with all uneliminated nonroot
circles, with the additional root loop and one remaining
two-cell. Substitution leaves the root exponent equal to
$1$, because the outer word $W_h$ has exponent zero in
every child. This conclusion holds for any chosen lifts of
the target coefficients.

In either case the coefficient group for that remaining
one-pair presentation is the torsion-free free product of
$H$ with a free group. Forester--Rourke's unimodular
adjunction theorem, in the exact form checked in A.18,
gives asphericity and an injective coefficient map for every
$Y_N$; see
\cite[Theorem~1 and Proposition~1]{CpathsForesterRourke2005Adjunction}
and \cite[\S\S1--3]{CpathsForesterRourkeDiagrams}.
It applies to $N=0$ directly as well.

Take the actual nested CW union $Y=\bigcup_{N<\omega}Y_N$.
Every map of a sphere of dimension at least two has image
in a finite subcomplex and thus lies in some aspherical
$Y_N$, where it is nullhomotopic. Hence $Y$ is aspherical.
If a loop from $B$ became nullhomotopic in $Y$, the image of
its nullhomotopy would lie in a finite prefix, where
coefficient injection implies that the original element
of $H$ is trivial. This proves $H\hookrightarrow P$ without
requiring injectivity between consecutive prefix groups.
The relative cellular boundary is the identity computed
above. It gives the ordinary integral homology isomorphism
in (A.27.1), and the dimensions follow from the actual cell
attachments.

To verify normal generation, quotient the displayed
presentation by the normal closure of $H$. In the chain
case the first equation gives $x_0=1$, and all remaining
equations then give $x_n=1$. In the tree case its root
equation has $W_1=1$, hence $x_o=1$; every other coefficient
$x_oe_n^{-1}$ is then also $1$, so all other variables are
$1$. Thus the quotient is trivial in both cases. Finally,
any homomorphism $P\to E$ whose restriction to $K$ is $q$
would assign solutions in
$E$ to the displayed variables, which we have proved
impossible.

For the existential-closedness assertions, a nonidentity
mixed word $w\in E*\langle t\rangle$ satisfies $w(t)\ne1$
in that actual free-product overgroup. Existential
closedness realizes this single finite inequation in $E$.
For a torsion-free coefficient group the free product is
torsion-free as well, so the same argument applies within
that class. This proves property~1 in both cases.

Lastly, given an epimorphism $G\twoheadrightarrow Q$ and
a countable equation system over $Q$, choose lifts of all
occurring coefficients. The lifted words have exactly the
same exponent matrix. Any solution in $G$ maps to a
solution in $Q$. This proves the stated permanence and the
quotient obstruction. It does not assert that every
countable target containing a prescribed perfect coefficient
group has a quotient in one of the listed classes.
\end{proof}
